% Appendix A: notes on the toll registers.
\chapter{The Toll Registers}
\label{app:registers}

This appendix describes the conventions used when transcribing the registers listed in Table~\ref{tab:registers}.

\section{Transcription}

Entries were transcribed in full, preserving the original spelling of names and goods. Dates given by saints' days were converted to the Julian calendar using the standard tables. Where an entry is damaged, the legible portion is recorded and the remainder marked as missing; no entries were reconstructed by conjecture.

\section{Classification of Goods}

Goods were assigned to eleven categories: salt, herring, grain, beer, cloth, wax, furs, metals, timber, wine and miscellaneous. Mixed cargoes were assigned to the category named first in the entry, on the assumption, supported by the tariffs, that toll keepers listed the most valuable item first.

\section{Estimating Weights}

For salt, one barrel is taken as 140 kilograms and one wagon as ten barrels. For other goods, the conversion factors follow the tariffs of the respective stations. Entries recording only a payment were converted using
\begin{equation}
  w = \frac{p}{r_g},
\end{equation}
where $p$ is the toll paid in pennies and $r_g$ is the station's rate per hundredweight for goods of category $g$.
